\section{Completion of valued fields}

\begin{definition}[Cauchy sequence]
    Let \(K\) be a field equipped with a metric
    \[
        d:K\times K\longrightarrow \mathbb{R}_{\geq 0}.
    \]
    A sequence
    \[
        (a_n)_{n\geq 1}=(a_1,a_2,\dots,a_n,\dots)
    \]
    of elements of \(K\) is called a Cauchy sequence if for every
    \[
        \varepsilon>0,
    \]
    there exists an integer \(N\geq 1\) such that
    \[
        d(a_m,a_n)<\varepsilon
        \qquad
        \text{for all } m,n\geq N.
    \]
\end{definition}

\begin{definition}[Convergent sequence]
    Let \(K\) be a field equipped with a metric \(d\). A sequence
    \[
        (a_n)_{n\geq 1}
    \]
    in \(K\) is said to converge to an element \(a\in K\) if
    \[
        \lim_{n\to\infty} d(a_n,a)=0.
    \]
    In this case, we write
    \[
        \lim_{n\to\infty} a_n=a.
    \]

    If the metric \(d\) is induced by a multiplicative valuation \(\varphi\), namely
    \[
        d(x,y)=\varphi(x-y),
    \]
    then the above condition is equivalently
    \[
        \lim_{n\to\infty} \varphi(a_n-a)=0.
    \]
\end{definition}

\begin{definition}[Complete valued field]
    Let \(K\) be a field equipped with a metric \(d\), or equivalently with a
    valuation inducing \(d\). We say that \(K\) is complete if every Cauchy
    sequence in \(K\) converges to an element of \(K\).

    If \(K\) is complete with respect to the metric induced by a valuation
    \(\varphi\), then \(K\) is also called a complete valued field.
\end{definition}

\begin{definition}[Complete valuation ring]
    Let \((K,\varphi)\) be a valued field, and let
    \[
        R_{\varphi}=\{x\in K\mid \varphi(x)\leq 1\}
    \]
    be its valuation ring in the multiplicative convention.
    If \(K\) is complete with respect to the metric induced by \(\varphi\), then
    \(R_{\varphi}\) is called a complete valuation ring.
\end{definition}

\begin{definition}[Ring of Cauchy sequences]
    Let \((K,\varphi)\) be a valued field, and let
    \[
        d(a,b)=\varphi(a-b)
    \]
    be the metric induced by \(\varphi\). Denote by \(\mathcal{C}\) the set of all
    Cauchy sequences in \(K\):
    \[
        \mathcal{C}
        :=
        \left\{
            (a_n)_{n\geq 1}
            \ \middle|\
            (a_n)_{n\geq 1}\text{ is a Cauchy sequence in }K
        \right\}.
    \]
    For \((a_n),(b_n)\in \mathcal{C}\), define
    \[
        (a_n)+(b_n):=(a_n+b_n),
    \]
    and
    \[
        (a_n)(b_n):=(a_nb_n).
    \]
    With these operations, \(\mathcal{C}\) is a commutative ring.
\end{definition}

\begin{definition}[Null sequences]
    Let
    \[
        \mathcal{I}
        :=
        \left\{
            (a_n)\in \mathcal{C}
            \ \middle|\
            \lim_{n\to\infty} a_n=0
        \right\}.
    \]
    Equivalently,
    \[
        (a_n)\in \mathcal{I}
        \Longleftrightarrow
        \lim_{n\to\infty}\varphi(a_n)=0.
    \]
    The elements of \(\mathcal{I}\) are called null sequences.
\end{definition}

\begin{proposition}
    The subset \(\mathcal{I}\) is a maximal ideal of \(\mathcal{C}\). Consequently,
    the quotient
    \[
        \widehat{K}:=\mathcal{C}/\mathcal{I}
    \]
    is a field.
\end{proposition}

\begin{proof}
    It is straightforward to check that \(\mathcal{I}\) is an ideal of
    \(\mathcal{C}\).

    To show that \(\mathcal{I}\) is maximal, let
    \[
        (a_n)\in \mathcal{C}\setminus \mathcal{I}.
    \]
    Then \((a_n)\) does not converge to \(0\). Since \((a_n)\) is Cauchy, there exist
    \(N\geq 1\) and \(\varepsilon>0\) such that
    \[
        \varphi(a_n)\geq \varepsilon
        \qquad
        \text{for all } n\geq N.
    \]
    Hence \(a_n\neq 0\) for all sufficiently large \(n\). Define a sequence
    \[
        b_n=
        \begin{cases}
            0, & n<N,\\
            a_n^{-1}, & n\geq N.
        \end{cases}
    \]
    Then \((b_n)\) is again a Cauchy sequence in \(K\), and
    \[
        (a_n)(b_n)-1\in \mathcal{I}.
    \]
    Therefore the class of \((a_n)\) in \(\mathcal{C}/\mathcal{I}\) is invertible.
    Thus every nonzero element of \(\mathcal{C}/\mathcal{I}\) is invertible, so
    \(\mathcal{C}/\mathcal{I}\) is a field. Hence \(\mathcal{I}\) is maximal.
\end{proof}

\begin{definition}[Completion of a valued field]
    The field
    \[
        \widehat{K}:=\mathcal{C}/\mathcal{I}
    \]
    is called the completion of \(K\) with respect to the valuation \(\varphi\).
\end{definition}

\begin{proposition}[Extension of the valuation]
    Let \((a_n)\in \mathcal{C}\). Then the sequence
    \[
        \bigl(\varphi(a_n)\bigr)_{n\geq 1}
    \]
    converges in \(\mathbb{R}_{\geq 0}\).

    Moreover, if
    \[
        (a_n)-(b_n)\in \mathcal{I},
    \]
    then
    \[
        \lim_{n\to\infty}\varphi(a_n)
        =
        \lim_{n\to\infty}\varphi(b_n).
    \]
    Hence the rule
    \[
        \widehat{\varphi}\bigl((a_n)+\mathcal{I}\bigr)
        :=
        \lim_{n\to\infty}\varphi(a_n)
    \]
    defines a well-defined multiplicative valuation on \(\widehat{K}\).
\end{proposition}

\begin{proof}
    Since \((a_n)\) is Cauchy in \(K\), for every \(\varepsilon>0\) there exists
    \(N\geq 1\) such that
    \[
        \varphi(a_n-a_m)<\varepsilon
        \qquad
        \text{for all } m,n\geq N.
    \]
    By the reverse triangle inequality,
    \[
        \left|\varphi(a_n)-\varphi(a_m)\right|
        \leq
        \varphi(a_n-a_m).
    \]
    Therefore \(\bigl(\varphi(a_n)\bigr)_{n\geq 1}\) is a Cauchy sequence in
    \(\mathbb{R}_{\geq 0}\). Since \(\mathbb{R}\) is complete, the limit
    \[
        \lim_{n\to\infty}\varphi(a_n)
    \]
    exists.

    Now assume that
    \[
        (a_n)-(b_n)\in \mathcal{I}.
    \]
    Then
    \[
        \lim_{n\to\infty}\varphi(a_n-b_n)=0.
    \]
    Again by the reverse triangle inequality,
    \[
        \left|\varphi(a_n)-\varphi(b_n)\right|
        \leq
        \varphi(a_n-b_n).
    \]
    Taking limits gives
    \[
        \lim_{n\to\infty}\varphi(a_n)
        =
        \lim_{n\to\infty}\varphi(b_n).
    \]
    Thus \(\widehat{\varphi}\) is well-defined.

    Finally, for \((a_n),(b_n)\in \mathcal{C}\), we have
    \[
        \widehat{\varphi}
        \left(
            \bigl((a_n)+\mathcal{I}\bigr)
            \bigl((b_n)+\mathcal{I}\bigr)
        \right)
        =
        \lim_{n\to\infty}\varphi(a_nb_n)
        =
        \lim_{n\to\infty}\varphi(a_n)\varphi(b_n),
    \]
    hence
    \[
        \widehat{\varphi}
        \left(
            \bigl((a_n)+\mathcal{I}\bigr)
            \bigl((b_n)+\mathcal{I}\bigr)
        \right)
        =
        \widehat{\varphi}\bigl((a_n)+\mathcal{I}\bigr)
        \widehat{\varphi}\bigl((b_n)+\mathcal{I}\bigr).
    \]
    Similarly, the triangle inequality for \(\varphi\) passes to the limit. Hence
    \(\widehat{\varphi}\) is a multiplicative valuation on \(\widehat{K}\).
\end{proof}

\begin{proposition}
    There is a natural embedding of valued fields
    \[
        K\hookrightarrow \widehat{K}.
    \]
\end{proposition}

\begin{proof}
    Define
    \[
        \iota:K\longrightarrow \widehat{K}
    \]
    by sending \(a\in K\) to the class of the constant sequence:
    \[
        \iota(a):=(a,a,a,\dots)+\mathcal{I}.
    \]
    This is a field homomorphism. If \(\iota(a)=0\), then the constant sequence
    \((a,a,a,\dots)\) belongs to \(\mathcal{I}\), so
    \[
        \lim_{n\to\infty}\varphi(a)=0.
    \]
    Hence \(\varphi(a)=0\), and therefore \(a=0\). Thus \(\iota\) is injective.

    Moreover,
    \[
        \widehat{\varphi}(\iota(a))
        =
        \widehat{\varphi}\bigl((a,a,a,\dots)+\mathcal{I}\bigr)
        =
        \varphi(a),
    \]
    so the embedding preserves the valuation.
\end{proof}

\begin{theorem}[Completion of a valued field]
    Let \((K,\varphi)\) be a valued field. Then there exists a complete valued field
    \[
        (\widehat{K},\widehat{\varphi})
    \]
    together with a valuation-preserving embedding
    \[
        K\hookrightarrow \widehat{K}
    \]
    such that the image of \(K\) is dense in \(\widehat{K}\).

    Moreover, \(\widehat{K}\) is unique up to a unique \(K\)-isomorphism of valued
    fields.
\end{theorem}

\begin{proof}
    The field \(\widehat{K}=\mathcal{C}/\mathcal{I}\), together with the valuation
    \(\widehat{\varphi}\), is complete. Indeed, every Cauchy sequence in
    \(\widehat{K}\) can be represented by a suitable diagonal sequence of Cauchy
    sequences in \(K\), and this diagonal sequence gives its limit in
    \(\widehat{K}\).

    The image of \(K\) is dense in \(\widehat{K}\). If
    \[
        x=(a_n)+\mathcal{I}\in \widehat{K},
    \]
    then the sequence of constant classes
    \[
        \iota(a_n)\in \widehat{K}
    \]
    converges to \(x\).

    Finally, if \(L\) is another complete valued field containing \(K\) densely and
    preserving the valuation, then every Cauchy sequence in \(K\) has a unique limit
    in \(L\). Thus the rule
    \[
        (a_n)+\mathcal{I}
        \longmapsto
        \lim_{n\to\infty}a_n
    \]
    defines a unique \(K\)-isomorphism
    \[
        \widehat{K}\cong L
    \]
    of valued fields. Hence the completion is unique up to \(K\)-isomorphism of
    valued fields.
\end{proof}  
